

\chapter{[AUU] Anamnese und Patientengespräch}
\label{chp:AUU}

\emph{Maintainer: Sebastian Gabriel}

\minitoc

Hier lernen Sie wie Sie die bereits gelernten theoretischen Grundlagen in
der Praxis einsetzen, wie und in welcher Reihenfolge Sie im Einsatz
vorgehen, wie Sie Dringlichkeiten einsch\"atzen und Verdachtsdiagnosen
erarbeiten können und wie Sie dabei dann auch noch den Überblick
behalten können.


%========Beginn Abschnitt Emhofer============================================

\section{{\quotedblbase}Hilfe, ein Patient!{\textquotedblleft} - Was
mach{\textquotesingle} ich jetzt?}

\begin{wrapfigure}{r}{4cm}
\includegraphics[width=\linewidth]{./images/teaser/follow-me-kontur.jpg}
\end{wrapfigure}Gerade zu Beginn der T\"atigkeit als medizinische Fachkraft
fühlt man sich beim Eintreffen am Notfallort durch die große Menge an Sinneseindrücken und
Informationen überlastet. Um Sie sicher durch diese Informationsflut zu
leiten, wollen wir Ihnen in diesem Abschnitt ein paar Anhaltspunkte
pr\"asentieren, welche den Einsatzablauf strukturieren und Ihnen und somit auch
dem Patient Sicherheit geben.

\Layout{0}{Vorgehen am Notfallort:}{
\begin{itemize}
  \item Bevor es los geht: Szeneüberblick und Selbstschutzmaßnahmen
  \item 1. Schritt: \Es{Ersteinsch\"atzung des Patienten}
  \item 2. Schritt: \Es{Erstmaßnahmen}
  \item 3. Schritt: \Es{Durchatmen, nachdenken, planen}
\end{itemize}
}{
\begin{itemize}
  \item Szeneüberblick und Selbstschutzmaßnahmen
  \item Ersteinsch\"atzung des Patienten
  \item Erstmaßnahmen
  \item Durchatmen, nachdenken, planen
\end{itemize}
}{}{}{}

\dictum[Dickie, in: Samuel Shem: House of God.
\cite{HouseOfGodDe}]{Alles in Ordnung, auch wenn Sie in Panik geraten sind und sich beschissen fühlen.
Ich kenne das. Es ist furchtbar, und es wird nicht das letzte Mal sein.
Aber vergessen Sie nicht, was Sie gesehen haben. Regel Nr. 3: Bei Herzstillstand
zuerst den eigenen Puls fühlen.}

W\"ahrend die ersten zwei Schritte ein gezieltes Vorgehen erfordern, bietet der
dritte Schritt Spielraum für die individuelle Untersuchung und Befragung des
Patienten entsprechend den Symptomen. Im Folgenden stellen wir Ihnen die Schritte
im Einzelnen vor.

% \section{Die Anamnese erz\"ahlt uns die
% Geschichte des Patienten}
\index{Anamnese}

\section{Kommunikationsregeln}
\label{topic:kommunikationsregeln}

\begin{wrapfigure}{r}{4cm}
\includegraphics[width=\linewidth]{./images/teaser/imgp0482.jpg}
\end{wrapfigure}
Der Erfolg des Anamnesegespräches ist für das weitere Vorgehen 
entscheidend. Oft erhält man wichtige Information zum Zustand des Patienten
nicht durch die gemessenen Werte sondern durch die richtige Befragung! Der
Erfolg des Gespräches hängt dabei einerseits von Geschick des Gesprächführenden,
andererseits aber auch von der Bereitschaft des Patienten (oder dessen
Umgebung!) ab. 

Im Folgenden möchten wir Ihnen einige \E{Grundregeln} und wertvolle Tipps
geben, wie Sie Ihr Patientengespräch so gestalten können, dass Sie die
gewünschten Informationen erhalten und ein Vertrauensverhältnis mit dem
Patienten aufbauen können:

\begin{itemize}
    \item Stellen Sie sich mit Ihrem Namen und Ihrer Funktion vor und - so es
    der Notfall zulässt - geben Sie dem Patienten die Hand!
    \item Begeben Sie sich auf die gleiche Ebene (auf \E{Augenhöhe}) mit dem
    Patienten!
    \item Sprechen Sie laut und deutlich!\footnote{Evtl. benötigt der
    Patient Hörhilfen. Diese müssen evtl. erst eingeschalten werden.}
    \item Hören Sie "`aktiv zu"'! Die Antworten sollen bei Ihnen "`ankommen"'
    und nicht "`verloren gehen"'. Wichtige Informationen werden oft nur einmal
    berichtet, der Patient nimmt im weiteren Verlauf an, dass er sie bereits
    gesagt und das Fachpersonal sie verstanden hat.
    
    Negativbeispiel: \begin{inparadesc}
                     \item[Sanitäter:] \Speech{"`Haben Sie Schmerzen?"'}
                     \item[Patient:] \Speech{"`Nein."'}
                     \item[Sanitäter:] \Speech{"`Wie stark ist der Schmerz?"'}
                     \end{inparadesc}
    \item Versuchen Sie zu verstehen was
    der Patient \E{meint}! Wiederholen Sie wichtige Informationen um sicher zu
    gehen dass Sie es richtig verstanden haben.
    \item Verwenden Sie eine leicht verständliche Sprache und vermeiden Sie Fachausdrücke!
    \item Stellen Sie immer nur \E{eine Frage gleichzeitig}.
    \item Informieren Sie den Patient über Ihr Vorgehen!
    \item Versuchen Sie Fragen des Patienten nach bestem Wissen zu
    beantworten, aber erfinden Sie nichts! Lügen Sie den Patienten nicht an!
    \item Treffen Sie keine Aussagen oder Vorhersagen die ausserhalb Ihres
    Kompetenz- und Zuständigkeitsbereiches liegen.
    \item Im Team sollten nicht alle auf den Patient einreden! Der
    Patient sollte wenn möglich nur einen Ansprechpartner
    haben!\label{topic:einer-redet}
\end{itemize}

\MarriedNG{Fragentypen}
{
Grundsätzlich unterscheidet man \Es{offene} und \Es{geschlossene} Fragen. 

\subparagraph{Offene Fragen} kann der Patient frei, mit seinen eignen Worten
beantworten. Je nach Patient kann man so mit wenig Aufwand sehr viele
Informationen gewinnen. Das Problem dabei kann jedoch sein, dass die
Informationen ungeordnet und mitunter auch in der Situation mäßig wichtig sind. 

\begin{quote}
  \Speech{"`Was haben Sie für ein Problem?"'}
  
  \Speech{"`Können Sie den Schmerz beschreiben?"'}
\end{quote}

\subparagraph{Geschlossene Fragen} sind sehr direkt und speziell, der Patient
kann i.d.R. nur mit \emph{"`ja"'}, \emph{"`nein"'} oder \emph{"`weiss nicht"'} antworten.
Sie machen daher eher nur dann Sinn, wenn man ganz spezielle Dinge bewusst
direkt abfragen kann oder muss. \Ca{Achtung!} Manche Patienten neigen dazu
grundsätzlich "`ja"' zu sagen, ohne dass sie die Frage wirklich verstanden
haben!

\begin{quote}
  \Speech{"`Haben Sie sich vorher angestrengt?"'}
  
  \Speech{"`Ist der Schmerz stechend?"'}
\end{quote}

\subparagraph{Suggestivfragen} sind ganz "`böse"': Hier
legt der Gesprächsführer den Patienten die Antworten "`in den Mund"'. 

\begin{quote}
  \Speech{"`Haben Sie \emph{eh keine} Schmerzen?"'}
  
  \Speech{"`Ist der Schmerz stechend?"'}\footnote{Ja, das ist die gleiche Frage wie
  das Beispiel bei der geschlossenen Frage. Man sieht, wie schnell es gehen kann,
  ohne es zu wollen eine Suggestivfrage zu stellen \ldots}
\end{quote}

Dieser Fragentyp
soll nicht -- oder nur im begründeten Ausnahmefall -- angewendet werden. 

}
{
% TODO VIT: Slide fehlt!
\begin{itemize}
    \item Offene Fragen.
    \item Geschlossene Fragem.
    \item Suggestivfragen (böse!).
\end{itemize}
}

\MarriedNG{Kinder}
{ Bei Notfällen mit Kindern kann die Kommunikation besonders herausfordernd
werden, da Kinder, die Angst haben, sich oft gegenüber Fremden verschließen. Die
oben genannten Regeln gelten daher bei Kindern umso mehr. Zusätzlich können
folgende Tipps hilfreich sein:

\begin{itemize}
  \item Kinder nehmen non-verbale Signale stärker wahr als Erwachsene.
  Versuchen Sie daher durch einen freundlichen Gesichtsausdruck und eine
  offene Körperhaltung Vertrauen zu bilden!
  \item Erklären Sie dem Kind (altersgerecht) warum Sie gerufen wurden!
  \item Sprechen Sie das Kind direkt und mit seinem Namen an! Vermeiden  Sie
  es über das Kind in der 3. Person zu sprechen, damit das Kind nicht das
  Gefühl  bekommt, dass hinter seinem Rücken getuschelt wird. Dies würde dem
  Kind unnötig Angst machen!
  \item Seien Sie immer ehrlich! Kinder merken sehr schnell, wenn Sie lügen.
  Lügen  Sie ein Kind an, so wird das Kind schnell das Vertrauen in Sie
  verlieren  und jede weitere Untersuchung verweigern!
  \item Geben Sie dem Kind ein Spielzeug in die Hand!
  \item Trennen Sie das Kind niemals von der Bezugsperson!
  \item Hat das Kind Angst vor einer bestimmten Untersuchung oder einer
  Maßnahme, zeigen Sie diese zuerst an der Mutter oder an einem Stofftier vor!
  \item Loben Sie das Kind nach der Untersuchung bzw. nach der Maßnahme!
  \item Da Kinder ihre Beschwerden nicht eindeutig beschreiben können, lassen
  Sie sich z.\,B. den Ort der Schmerzen zeigen. Auf die \E{Körperhaltung}
  besonders achten!
\end{itemize}
}{
{\SymbolText}
}


\section{Was muss ich erfragen? -- S-KAMEL}

\index{S-KAMEL}

Neben der körperlichen Untersuchung gibt das Gespr\"ach mit dem Patienten
wertvolle Hinweise auf das Notfallgeschehen. Diese Erhebung der Kranken- bzw.
Unfallgeschichte nennt man \T{Anamnese}. Eine gute Anamnese ist mindestens
genauso wichtig wie die Erhebung der Befunde! Leider wird in der Praxis oft der
Messwerterhebung eine größere Bedeutung beigemessen.

\Layout{0}{S-KAMEL}{
Auch das Anamnesegespräch soll strukturiert durchgeführt werden. Dafür stehen
mehrere Abfrageschemata zur Verfügungen, von welchen wir Ihnen ein in Österreich
weit verbreitetes vorstellen möchten: Das \T{S-KAMEL-Schema}. Nach einer
Einstiegsfrage, z.\,B. "`Was können wir für Sie tun?"', beinhaltet das
\Abk{S-KAMEL}-Schema folgende Punkte:

\begin{itemize}
    \item Symptome und Schmerzen
    \item Krankheiten
    \item Allergien
    \item Medikamente
    \item Ereignisse vor Notfalleintritt
    \item Letzte \ldots
\end{itemize}

Dieses Abfrageschema kann Ihnen im hektischen Notfall
Sicherheit geben. Durch das Abarbeiten der Fragen haben Sie
die Sicherheit, dass keine wichtigen Informationen vergessen
werden. Außerdem ist es für den Patienten vertrauensbildend,
da die gezielte Befragung professionell wirkt. Im Folgenden
müssen wir uns die einzelnen Punkte des Schemas genauer
ansehen.
}{
\begin{itemize}
    \item Symptome und Schmerzen
    \item Krankheiten
    \item Allergien
    \item Medikamente
    \item Ereignisse vor Notfalleintritt
    \item Letzte \ldots{}
\end{itemize}
}{}{}{}

\Fazit{S{\textquotesingle}KAMEL denkt an alles.}

\subsubsection{Symptome und Schmerzen}
\index{Schmerz}

Die Einstiegsfrage hat uns u.U. bereits auf das Hauptsymptom hingewiesen. Ist
dies nicht der Fall, so muss als erstes nach dem Hauptsymptom gefragt werden.
Fragen können z.\,B. sein:
\begin{itemize}
    \item \Speech{"`Welche Beschwerden haben Sie?"'}
    \item \Speech{"`Wo drückt der Schuh?"'}
\end{itemize}
Außerdem muss immer auch speziell nach Schmerzen gefragt werden!
\begin{itemize}
    \item \Speech{"`Haben Sie Schmerzen?"'}
\end{itemize}

\Layout{0}{Fragen über Schmerzen:}{
Gibt der Patient an Schmerzen zu haben, müssen wir weitere Informationen
erfragen, damit wir uns ein klares Bild über die Beschwerden machen können:
\index{Schmerz!Arten}

\begin{itemize}
    \item \Es{Lokalisation}: \Speech{"`Wo sind die
    Schmerzen?"'}
    \item \Es{Zeitdauer/Beginn}: \Speech{"`Seit wann haben Sie
    die Schmerzen? Wie lange haben Sie die Schmerzen schon?
    Haben Sie so einen Schmerz schon einmal gehabt? Haben die
    Schmerzen plötzlich oder langsam eingesetzt?"'} 
    \item
    \Es{Qualität}: \Speech{"`Wie ist der Schmerz? Bohrend,
    stechend, ziehend, krampfartig, auf- und
    abschwellend?"'}
    \item \Es{Ausstrahlung}: \Speech{"`Strahlt der Schmerz
    irgendwohin aus?"'} 
    \item \Es{Provokation}: \Speech{"`Was
    macht den Schmerz besser oder schlechter?"'} 
    \item \Es{Stärke}: \Speech{"`Wie stark ist der
    Schmerz?"'} \Z{In einigen Lehrbüchern wird empfohlen, den Schmerz anhand
    einer Skala von 1-10 beurteilen zu lassen, wir finden
    jedoch, dass auch eine solche Beurteilung recht subjektiv
    ist und daher nur mit Vorsicht verwendet werden sollte.}
    Setzen Sie die Aussage des Patienten immer in Relation zu
    seinem Gesichtsausdruck und zu seiner Körperhaltung!
    Einem Patient mit starken Schmerzen sieht man dies an!
\end{itemize}
}{
\begin{itemize}
  \item Wo?
  \item Ausstrahlung?
  \item Seit wann? Wie lange schon?
  \item Plötzlicher oder langsamer Beginn?
  \item Verlauf?
  \item Wie?
  \item Wie stark?
  \item Was provoziert den Schmerz?
\end{itemize}
}{}{}{}

\paragraph{Besonderheit: \T{Kolikartiger Schmerz}}~
\label{topic:kolik}
\index{Kolik!Schmerzqualität}
\index{Schmerz!kolikartiger}

\Definition{Ein kolikartiger Schmerz ist auf-
und abschwellend.}

Er wird oft als krampfhaft und sehr stark beschrieben. Er
kommt häufig bei Verstopfung der Gallen- und Harnwege durch
Steine vor (\See{topic:gallenkolik},
\ref{topic:nierenkolik}).

\subsubsection{Krankheiten}
Notfälle ereignen sich manchmal aufgrund der akuten Verschlechterung einer
bestehenden Vorerkrankung. Außerdem hilft die Beurteilung der Messwerte, wenn man
über die Grunderkrankungen des Patienten Bescheid weiß. Zum Beispiel kann die
Ursache eines hohen Blutdrucks eine bereits bestehende Erkrankung aber auch ein
Notfallgeschehen sein. Neben der Befragung kann auch ein alter Spitalsbericht
(Arztbrief) Auskunft über bestehende Vorerkrankungen liefern. 

Allgemeine Fragen,
wie z.\,B.
\begin{quote}
    \Speech{"`Leiden Sie unter irgendwelchen
    Krankheiten?"'}
    
    \Speech{"`Bestehen bei Ihnen Vorerkrankungen?"'}
\end{quote}
können anschließend durch symptom-spezifische Fragen ergänzt werden:
\begin{quote}
    \Speech{"`Hatten Sie schon einmal einen
    Herzinfarkt?"'}
    
    \Speech{"`Leiden Sie unter hohem Blutdruck?"'}
    
    \Speech{"`Sind Sie zuckerkrank?"'}
\end{quote}

\subsubsection{Allergien}
Auf die Frage nach möglichen Allergien wird häufig vergessen, ein
Informationsmangel kann aber fatale Auswirkungen haben. Fragen Sie daher immer
nach Allergien, in einer zweiten Frage sogar speziell nach Medikamentenallergien!
Damit können Sie für Not- oder Spitalsarzt wichtige Informationen erheben.
(Natürlich könnten das die Ärzte auch selbst fragen, aber denken Sie daran, dass
ein Notfallpatient bewusstlos werden kann und es deshalb wichtig ist, dass sie
diese Informationen einholen, solange der Patient noch ansprechbar ist!)

\subsubsection{Medikamente}
Die vom Patienten eingenommenen Medikamente lassen
Rückschlüsse auf seine Grunderkrankungen zu und sind im Falle
eines längeren Spitalsaufenthalts eine wertvolle Information.
Auch die Regelmäßigkeit oder Unregelmäßigkeit der Einnahme
kann Ihnen Anhaltspunkte für das Finden Ihrer
Verdachtsdiagnose bzw. der Umstände des Patienten liefern. Angaben zu
Medikamenten sind für die nachbehandelnde Einrichtung zu dokumentieren. 

\Layout{0}{Fragen zu Medikamenten:}{
\begin{quote}
    \emph{"`Welche Medikamente nehmen Sie?"'}
    
    \emph{"`Wogegen nehmen Sie diese Medikamente?"'}
    
    \emph{"`Nehmen Sie Ihre Medikamente regelmäßig?"'}
    
    \emph{"`Wann haben Sie Ihre Medikamente zuletzt"'}
    genommen?
\end{quote}
}{
\begin{itemize}
  \item Welche?
  \item Wogegen?
  \item Regelmäßige Einnahme?
  \item Wann zuletzt genommen?
\end{itemize}
}{}{}{}

\subsubsection{Ereignisse vor Notfalleintritt}
Oft ereignen sich vor einem Notfall Schlüsselereignisse, welche den Notfall
herbei führen oder diesen auslösen. Fragen Sie daher immer nach möglichen
Besonderheiten oder Tätigkeiten des Patienten vor dem Notfall:
\begin{quote}
    \emph{"`Was haben Sie vor Beginn der Symptome
    gemacht?"'}
    \emph{"`Haben Sie sich vor Beginn der Symptome
    besonders angestrengt?"'}
    u.v.m.
\end{quote}
In dieser Kategorie kann es auch hilfreich sein, wenn Sie Angehörige oder Zeugen befragen:
\begin{itemize}
    \item Hat der Patient vor dem Notfall etwas Eingenommen (Überdosis)?
    \item Hat der Patient gekrampft bevor er bewusstlos liegen geblieben ist?
    \item Ist der Patient vor dem Eintreten der Symptome gestürzt?
    \item u.v.m.
\end{itemize}

\subsubsection{Letzte ...}
Der letzte Punkt des Abfrageschemas bezieht sich auf letzte
Ereignisse, welche mit dem Notfall in Zusammenhang stehen
könnten. Damit Sie hier die richtigen Fragen stellen können,
ist ein fundiertes Fachwissen notwendig. Abhängig von den
Beschwerden können verschiedene Fragen zielführend sein.

\Layout{0}{Fragen zu Letzte\ldots}{
\begin{quote}
    \emph{"`Wann haben Sie zuletzt etwas
    gegessen/getrunken?"'} (Ist der Patient nüchtern?)
    
    \emph{"`Wann hatten Sie zuletzt Stuhlgang?"'}
    (weiterführende Frage: Wie hat der Stuhl ausgesehen?)
    
    \emph{"`Wann war die letzte Regelblutung?"'}
    (Besteht die Möglichkeit einer Schwangerschaft?)
    
    \emph{"`Wann waren Sie zuletzt beim Arzt bzw. im
    Spital?"'}
\end{quote}
}{
\begin{itemize}
  \item Letzte Mahlzeit?
  \item Letzter Stuhlgang?
  \item Letzte Regelblutung?
  \item Letzter Spitalsaufenthalt?
  \item Letzter Arztbesuch?
  \item Letzte Blutzuckerkontrolle?
  \item \ldots
\end{itemize}
}{}{}{}

\section{Der Arztbrief -- Ein besonderes Hilfsmittel für
die Anamnese}

\subsubsection{Beispiel: Patient Zapfl}







{\fontfamily{lmtt}\fontseries{lc}\selectfont




\begin{flushright}
 Krankenanstalt der Barmherzigkeit \\II. Med. Abteilung
 
 Hartlgasse 16a \\A-1240 Wien
 
+43\,555\,328\,745\,0\\KADBMED2@grissu.org

\end{flushright}

                                                                                               

\vspace{1.5em}
                                                                                                                            


 Peter Zapfl\par
Diesseitsweg 13a/12/5/34 \par
A-1147 Wien
\vspace{3em}

\begin{flushright}
Wien, am 25.07.2009
\end{flushright}

  Wir berichten über den Aufenthalt unseres Patienten Herrn Peter ZAPFL, geb.
  am 13.12.1936, welcher vom 30.03.2009 bis 17.04. 2009 an unserer Abteilung
  stationär aufgenommen war.

{\bfseries   Aus der Anamnese:}

\begin{quote}
  Kindheit: \MarginNo{\Z{AE: Entfernung des Blinddarmes; TE:
  Entfernung der Mandeln.}}AE, sowie TE, WK II: längerer Lazarettaufenthalt wg.
  Typhus.
  
  Seit 2003 sind ein bis dato nicht insulinpflichtiger Diabetes mellitus, sowie
  eine art. Hypertonie bekannt.
  
  2005 Aufenthalt an unserer Abteilung wg. eines im \MarginNo{\Z{CT:
  Computertomographie.}}CT nachgewiesenen rechtshirnigen ischämischen Insultes,
  anschl. Rehabaufenthalt in Bad Schallerbach.
  
  Zuletzt lag Pat vom 21.11.2008 bis 23.12.2008 wegen eines nicht
  insulinpflichtigen, entgleisten Diabetes mell., bei Z. n ischäm. Zerebralem.
  Insult 2005 mit \MarginNo{\Z{HSZ: Halbseitenzeichen.}}HSZ li und Art.
  Hypertonie an unserer Abteilung.
  
  Zwischenzeitlich keine Auffälligkeiten.
  
  Alkohol.: verneint. Nikotin: verneint. Koffein: fallweise. Harn, Stuhl: unauff.
  Schlaf: unauff. Ven.: neg
\end{quote}


{\bfseries Bisherige Med.:} 

\begin{quote}
  \begin{tabular}{lll}
Renitec 10 mg &&1-1-1-
\\
Diabetex 850 mg &&1-1-1
\\
Diab Diät 12 BE &&2-1-4-2-2-1
\\
Thromboass &&1-0-0
\\
\end{tabular}
\end{quote}





 Die nunmehrige Aufnahme erfolgte mit dem ASB, der wegen zunehmender
 Verschlechterung des AZ des Pat. vom Wohnungsnachbarn berufen wurde.

{\bfseries Auffälliges im Aufnahmestatus:}

\begin{quote}
 Deutlich herabgesetzter \MarginNo{AZ: Allgemeinzustand\Z{, EZ:
 Ernährungszustand}.}AZ, noch ausreichender EZ, zeitlich und örtlich
 teilorientiert, motorische Unruhe. Aktive Beweglichkeit der re.
 \MarginNo{OE: Obere Extremität, UE: Untere Extremität.}OE und UE, beinbetonte
 spastische Halbseitenlähmung li mit gesteigerten Sehnenreflexen li bei Z.n. re-Hirn-Insult, mäßige Artikulationsstörungen.
 
 Haut trocken, \MarginNo{\Z{Turgor: vom Flüssigkeitsgehalt abhängiger
 Spannungszustand der Haut. \nocite{Pschy259}}}Turgor herabgesetzt, kein
 Dekubitus nachweisbar.
 \end{quote}

{\bfseries Befunde:}
\begin{quote}
Labor siehe Ablichtung.

\Margin{\Z{C/P: "`Cor-Pulmo"', Thoraxröntgen.}}C/P.: Keine wesentl Änderung
gegenüber Vorbefund 12.2008

Schädel CT: Keine wesentliche Änderung gegenüber Letztbefund vom 12.2008,
Feinbefund siehe siehe Ablichtung

Augen:. Erstgradige hypertone Veränderungen der
\MarginNo{\Z{Retina: Netzhaut.}}Retinagefässe, kein HW für
\MarginNo{\Z{Diabetische Retinopathie: Schädigung der Netzhautgefäße
im Rahmen des Diabetes mellitus.}}diabet. Retinopathie.

Neuro: Anamnst Z. n. Insult 2005 im Strombereich der A. cerebri ant., mit
beinbetonter spastischer Halbseitenlähmung li und gesteigerten
\MarginNo{\Z{BSR, TSR: Sehnenreflexe des Bizeps und Trizeps.}}BSR und TSR li
sowie stark gesteigerten \MarginNo{\Z{PSR: Patellarsehnenreflex.
Babinski-Reflex: krankhafter Reflex bei Schädigung des Rückenmarks oder
Gehirns.}}PSR li, Babinski li pos, Trömner u. Knips bds neg, kein Hinweis auf
Re-Insult, weiterhin physi\-ka\-lisch- thera\-peutische und logopädische
Maßnahmen empf. Der durchgeführte Schlucktest ergab keinen Hinweis auf eine relevante Schluckstörung.

Uro: Prost. rect. dig. Unauff.

EKG: \MarginNo{\Z{Sinusrhythmus, HF 74, Linkstyp.}}SR 74, LT,
altersentsprechender Kurvenverlauf.

RR bei mehrfachen Kontrollen gegen Ende des Aufenthaltes normoton, zuvor im
hypotonen Bereich.
\end{quote}




{\bfseries Epikrise:}

\begin{quote}
  Nach Behebung des bei Aufnahme bestehenden Flüssigkeitsdefizites mit
  geeigneter Infusionstherapie, besserte sich der AZ des Pat. zusehends. Mit in der Folge
  durchgeführten physi\-ka\-lisch- therap. Maßnahmen konnte der
  Mobilisisationsgrad des Pat. günstig beeinflusst werden.
  
Der Pat. wird mit bis zum selbständigen Gehen mit Rollator mobilisiert entlassen.
Die zu Aufenthaltsbeginn durchwegs im hypotonen Bereich gemessenen RR-Werte haben
sich in mehrfachen Kontrollen unter dem angegeben antihypertonen Regime
normalisiert. Ebenso hat sich die ursprünglich unbefriedigend gefundene
diabetische Stoffwechsellage, bei konsequenter Einhaltung des vorgeschlagenen
Regimes im therapeutischen Bereich stabilisiert.

  Im Einverständnis mit dem Pat. und dessen Tochter wurde ein Platz im
  Bernhard-Prischl-Heim zur Gewährleistung einer regelmäßigen und auch
  ausreichenden Flüssigkeitszufuhr, sowie einer gesicherten Diabetes- Diät,
  organisiert. Eine Kontrolle der vor etwa einem Jahr durchgeführten 
  \MarginNo{\Z{Carotissonographie: Ultraschall der
  Halsschalgader.}}Carotissonographie ist angezeigt. Ein entspr. Termin wurde am
  27.12.2009  12,30 Uhr an der Angiolog Ambulanz d. KA Mitschka-Stiftung
  fixiert. Eine Befund\-be\-sprech\-ung ist nach Terminvereinbarung an h.o. Int.
  Ambulanz möglich.
\end{quote}

\MarginNo{\Z{}}
{\bfseries Hauptdiagnose:}

\begin{quote}
\begin{tabular}{ll}
Hypohydratation & E86.0
\\
\end{tabular}
\end{quote}

{\bfseries Nebendiagnosen:}

\begin{quote}
\begin{tabular}{ll}
entgleister. Diabetes mell. (NIDDM) & E12.31, E12.40, E12.72 
\\
Art. Hypertonie & I10.1
\\
Z.n. Zerebralem Insult mit HSS li 2005& I63.3
\\
Spastische Hemiparese li, beinbetont&I69.3, G81.1
\\
\end{tabular}
\end{quote}



{\bfseries Therapie:}

\begin{quote}
\begin{tabular}{lll}
Renistad Tab. 5,0 mg && 1-1-1 
\\
Glucophage 850 mg Tab && 1-1-1 
\\
Diab. Diät 12 BE&&(2-1-4-2-2-1) 
\\
Thromboass Tab.&&1-0-0 
\\
tgl. Flktszufuhr nicht unter 2 l&&
\end{tabular}
\end{quote}




{\bfseries Kontrollen:}

\begin{quote}
 NFP, Harnbefund, Elektrolyte, RR, BZ, (Pat. führt seit
Jahren Selbstkontrollen durch und führt darüber Aufzeichnungen)) weiterhin neurolog. und
ophthalmolog. Ko
\end{quote}









Mit kollegialen Grüßen\vspace{2em}\\ Dr. Adolf Schrammel, Stationsarzt

\Hand{Schrammel}

vidit OA Univ.Doz. Dr. Alois Schremser
}

\vspace{2em}

\hrule












